%=====================================================================
%  A Survey on Robust Multi-Agent Reinforcement Learning
%  Target venue: ACM Computing Surveys (CSUR)  ->  acmart / acmsmall
%  Skeleton built 2026-06-17. See robust-marl-MASTER.md for the 146-paper corpus.
%
%  STRUCTURE: modular, one file per section (one owner per file, see CONTRIBUTING.md).
%  Each section file header lists: OWNER, TARGET LENGTH, MASTER #s, and the
%  mandatory closing two questions.
%
%  THESIS (every section must serve this):
%    Single-agent robust-RL perturbation dimensions (model/state/action/reward)
%    all appear in MARL; MARL ADDS three dimensions absent in single-agent RL --
%    untrustworthy teammates, attacked communication, curse of multiagency --
%    which current methods do not yet handle in a unified way.
%
%  DIFFERENTIATION vs CSUR'25 (Adversarial ML in MARL): that SoK covers only the
%  attack/defence axis. We add DRMG theory (§3), LLM-MAS (§8), offline (§9), and a
%  unified "perturbation-source" taxonomy.
%=====================================================================
\documentclass[acmsmall,screen,review,anonymous]{acmart}
% acmsmall = ACM journal (CSUR) format. Remove `review,anonymous` for camera-ready.

\usepackage{booktabs}
\usepackage{multirow}
\usepackage{makecell}
% NOTE: acmart already loads amsmath + math-symbol fonts. Do NOT add amssymb
% (it clashes with acmart: "\Bbbk already defined").
\usepackage[edges]{forest}   % for the taxonomy tree figure (fig/taxonomy.tex)
\usepackage{xcolor}
\usepackage{pifont}          % \ding marks (amssymb is excluded above)
% Coverage marks for comparison tables: full / partial / (none = ``--'').
\providecommand{\checkmark}{\ding{51}}
\newcommand{\tcheck}{\ensuremath{\circ}}   % partial / peripheral coverage

%--- CSUR boilerplate (fill at submission) ---------------------------
\acmJournal{CSUR}
\acmVolume{0}\acmNumber{0}\acmArticle{0}\acmYear{2026}\acmMonth{0}
\citestyle{acmnumeric}

\begin{document}

\title{A Survey on Robust Multi-Agent Reinforcement Learning}
% Alt titles to consider:
%  - "Robust Multi-Agent Reinforcement Learning: A Survey of Perturbations, Methods, and Open Problems"
%  - "Beyond Single-Agent Robustness: A Survey of Robust Multi-Agent RL"

\author{Sihong He}
\affiliation{\institution{TODO}\country{USA}}
\email{sihonghe.ai@gmail.com}
\author{TODO Coauthors}
\affiliation{\institution{TODO}\country{TODO}}

\renewcommand{\shortauthors}{He et al.}

\begin{abstract}
% [Lead] ~200 words. (1) MARL is being deployed in safety-critical settings;
% (2) robustness is essential but MARL robustness is NOT a trivial extension of
% single-agent robust RL; (3) we organize the field by perturbation source and
% surface three MARL-specific dimensions; (4) we cover N=146 works incl. 2025-2026
% theory/LLM-MAS/offline that prior surveys miss; (5) we distill open challenges.
TODO.
\end{abstract}

\begin{CCSXML}
% TODO at submission. e.g. Computing methodologies~Multi-agent reinforcement learning
\end{CCSXML}
\ccsdesc[500]{Computing methodologies~Multi-agent systems}
\ccsdesc[500]{Computing methodologies~Reinforcement learning}

\keywords{multi-agent reinforcement learning, robustness, distributionally robust
Markov games, adversarial attacks, Byzantine fault tolerance, survey}

\maketitle

%======================  PAGE BUDGET (~36 pp body)  ===================
%  Abstract+Intro .............. 2.5   (01)
%  Background+Taxonomy ......... 3.0   (02)
%  §3 Model/Env uncertainty .... 5.0   (03)   <- deepest (theory)
%  §4 State/Obs perturbation ... 3.0   (04)
%  §5 Adversarial atk & train .. 3.5   (05)
%  §6 Communication ............ 2.5   (06)
%  §7 Teammates/Byzantine ...... 3.0   (07)
%  §8 LLM multi-agent safety ... 2.0   (08)   <- differentiation highlight
%  §9 Offline / dist. shift .... 1.0   (09)
%  §10 Benchmarks .............. 2.0   (10)
%  §11 Applications ............ 3.0   (11)
%  §12 Challenges & Future ..... 3.0   (12)   <- the payoff section
%  §13 Conclusion .............. 1.0   (13)
%======================================================================

%=====================================================================
% §1 INTRODUCTION                       OWNER: Lead   |  TARGET: ~2.5 pp (w/ abstract)
%=====================================================================
\section{Introduction}
\label{sec:intro}

% --- Outline (write as flowing prose, not bullets) ---
% (a) Motivation: MARL successes (SMAC, traffic, power grid, autonomous driving,
%     warehouse, LLM multi-agent) -> increasingly safety-critical deployment.
% (b) Problem: real deployments face uncertainty/attacks; non-robust MARL fails.
% (c) Gap / THESIS: robust MARL is NOT a trivial lift of single-agent robust RL.
%     Single-agent dims (model/state/action/reward) all recur, BUT MARL adds
%     three new dimensions -- untrustworthy teammates, attacked communication,
%     curse of multiagency.  <-- state this explicitly; it is the spine.
% (d) Differentiation vs prior surveys (see Sec.~\ref{sec:background} related work):
%     single-agent robust-RL surveys omit multi-agent dims; general MARL surveys
%     are not robustness-focused; CSUR'25 SoK covers only attack/defence and omits
%     DRMG theory, LLM-MAS, offline.  We give a unified perturbation-source taxonomy.
% (e) Contributions (bullet list):
%     1) unified taxonomy (Fig.~\ref{fig:taxonomy}) organizing 146 works 2017--2026;
%     2) first survey to integrate DRMG theory + adversarial + communication +
%        Byzantine + LLM-MAS + offline under one frame;
%     3) per-dimension comparison of assumptions/methods/guarantees;
%     4) distilled open challenges (Sec.~\ref{sec:challenges}).
% (f) Paper organization paragraph + forward-ref every section.

TODO prose.

%=====================================================================
%  Taxonomy figure (signature figure, place right after Introduction).
%  Built around the THESIS: shared-with-single-agent dims vs MARL-specific dims.
%  Uses forest package (\usepackage[edges]{forest} in main.tex).
%  This is a DRAFT layout -- tune colors/spacing at polish time.
%=====================================================================
\begin{figure*}[t]
\centering
\footnotesize
\begin{forest}
  for tree={
    grow=east, parent anchor=east, child anchor=west,
    draw, rounded corners, align=left, edge={-latex},
    l sep=10mm, s sep=2mm, anchor=west,
  }
  [Robust\\MARL, fill=gray!15
    [\textbf{I. Perturbation source}\\(shared w/ single-agent RL), fill=blue!8
      [{Environment / model uncertainty (\S3)\\\,DRMG, sample complexity, minimax}]
      [State / observation perturbation (\S4)]
      [Action perturbation (\S5)]
      [Reward poisoning (\S5)]
    ]
    [{\textbf{II. MARL-specific dimensions}\\(absent in single-agent RL; \emph{core of this survey})}, fill=red!8
      [Communication attacks \& noise (\S6)]
      [Untrustworthy teammates /\\Byzantine / fault tolerance (\S7)]
      [{Curse of multiagency\\(theory, \S3)}]
    ]
    [\textbf{III. Emerging frontiers}, fill=green!8
      [LLM multi-agent safety (\S8)]
      [Offline / distribution shift (\S9)]
    ]
    [\textbf{IV. Cross-cutting}, fill=yellow!12
      [Benchmarks \& evaluation (\S10)]
      [Applications (\S11)]
    ]
  ]
\end{forest}
\caption{Taxonomy of robust MARL. Branch~I collects perturbation dimensions
inherited from single-agent robust RL; Branch~II collects dimensions that are
\emph{specific to the multi-agent setting} and constitute the core argument of
this survey. A secondary methodological axis (distributionally robust
optimization, adversarial training, certified robustness, regularization) cuts
across all branches and is indicated per work in the section tables.}
\label{fig:taxonomy}
\end{figure*}
   % signature taxonomy figure referenced throughout

% Closing of intro: 1 short paragraph naming the three MARL-specific dimensions
% that recur as a refrain in every subsequent section.

%=====================================================================
% §2 BACKGROUND & RELATED WORK          OWNER: Lead (+ Theory owner for 2.2)
% TARGET: ~3 pp
% Goal: fix notation/terminology ONCE here so all later sections are consistent.
% DRAFT v1 (auto-drafted). Prose is synthesised, not enumerative; every \cite key
% resolves in refs.bib. Math uses amsmath only (acmart loads no amssymb).
%=====================================================================
\section{Background and Scope}
\label{sec:background}

This section fixes the notation and terminology used throughout the survey and
makes our organizing principle precise. We first recall the multi-agent
reinforcement learning (MARL) model (\S\ref{sec:bg-marl}), then disentangle the
overloaded vocabulary of \emph{robustness} and define the worst-case objects that
recur in every later section (\S\ref{sec:bg-robust}). We then state the
perturbation-source taxonomy that structures the survey (\S\ref{sec:bg-taxonomy})
and position our scope against prior surveys (\S\ref{sec:bg-related}).

\subsection{Multi-Agent RL Preliminaries}
\label{sec:bg-marl}

We model multi-agent interaction as a \emph{Markov game} (stochastic game)
\cite{littman1994markov}, the $n$-agent generalization of the Markov decision
process, written as the tuple
\[
\mathcal{G} = \bigl(\mathcal{N},\,\mathcal{S},\,\{\mathcal{A}_i\}_{i\in\mathcal{N}},
\,P,\,\{r_i\}_{i\in\mathcal{N}},\,\gamma\bigr),
\]
where $\mathcal{N}=\{1,\dots,n\}$ is the set of agents, $\mathcal{S}$ the state
space, $\mathcal{A}_i$ agent $i$'s action space, and
$\mathcal{A}=\mathcal{A}_1\times\cdots\times\mathcal{A}_n$ the \emph{joint} action
space. A joint action is $\mathbf{a}=(a_1,\dots,a_n)$, and $\mathbf{a}_{-i}$
denotes the actions of all agents other than $i$. The transition kernel
$P(s'\mid s,\mathbf{a})$ depends on the joint action; each agent receives a reward
$r_i(s,\mathbf{a})$, and $\gamma\in[0,1)$ discounts the future. Under partial
observability (a Markov game played over a Dec-POMDP), agent $i$ acts on a local
observation $o_i$ drawn from an observation model rather than the full state $s$.
Each agent follows a policy $\pi_i$ (history- or observation-conditioned), and the
\emph{joint policy} $\boldsymbol{\pi}=(\pi_1,\dots,\pi_n)$ induces the value
$V_i^{\boldsymbol{\pi}}(s)=\mathbb{E}\!\left[\sum_{t\ge0}\gamma^t
r_i(s_t,\mathbf{a}_t)\mid s_0=s\right]$. Table~\ref{tab:notation} collects this
notation.

\paragraph{Cooperation structure.} Three regimes recur. In the
\emph{cooperative} (team) setting all agents share a single reward
$r_i\equiv r$ and maximize a common return; value factorization
(e.g., QMIX \cite{rashid2018qmix}) and centralized-critic policy-gradient methods
(MADDPG \cite{lowe2017multi}, MAPPO \cite{yu2022surprising}) dominate here. In the
\emph{competitive} (zero-sum) setting agents have opposed rewards, and the
relevant solution concept is a (minimax) equilibrium. The \emph{mixed}
(general-sum) setting interpolates between the two. Most methods are trained under
\emph{centralized training with decentralized execution} (CTDE): a learner may use
global information at training time, but each agent must act on its own
observation at deployment---a gap that, as we will see, is exactly where many
robustness failures originate. We refer the reader to general MARL surveys
\cite{zhang2021multi,gronauer2022multi,jin2025comprehensive} for algorithmic
breadth; our focus is orthogonal.

\paragraph{Solution concepts.} Whereas single-agent RL seeks a single optimal
policy, MARL targets an \emph{equilibrium}: a joint policy from which no agent (or
coalition) benefits by unilateral deviation (Nash, correlated, or coarse
correlated equilibria), or, in the cooperative case, a team-optimal joint policy.
This shift matters for robustness: a perturbation does not merely degrade one
policy, it can move the entire equilibrium, and worst-case reasoning must be
carried out over joint behavior rather than a single agent's.

% Notation table for the survey. Referenced as Table~\ref{tab:notation} in §2.
\begin{table}[t]\centering\footnotesize
\caption{Notation used throughout the survey.}
\label{tab:notation}
\begin{tabular}{@{}ll@{\qquad}ll@{}}
\toprule
Symbol & Meaning & Symbol & Meaning \\
\midrule
$\mathcal{N},\,n$ & agent set, \# agents & $\pi_i,\,\boldsymbol{\pi}$ & policy of $i$, joint policy \\
$\mathcal{S},\,s$ & state space, state & $V_i^{\boldsymbol{\pi}}$ & value of $i$ under $\boldsymbol{\pi}$ \\
$\mathcal{A}_i,\,a_i$ & action space / action of $i$ & $\gamma$ & discount factor \\
$\mathcal{A},\,\mathbf{a}$ & joint action space / action & $\mathcal{U}$ & uncertainty set \\
$\mathbf{a}_{-i}$ & actions of all agents but $i$ & $\underline{V}$ & robust (worst-case) value \\
$o_i$ & local observation of $i$ & DRMG & distributionally robust Markov game \\
$P$ & transition kernel $P(s'\!\mid\! s,\mathbf{a})$ & CTDE & centralized train, decentral.\ exec. \\
$r_i,\,r$ & reward of $i$ / shared reward & Dec-POMDP & decentralized partially obs.\ MDP \\
\bottomrule
\end{tabular}
\end{table}


\subsection{What ``Robust'' Means: A Terminology Map}
\label{sec:bg-robust}

The words \emph{robust}, \emph{resilient}, \emph{safe}, and \emph{secure} are used
inconsistently across the literature. We adopt the following working
distinctions. \textbf{Robustness} is worst-case performance over a set of
admissible deviations from the nominal model or inputs---the central object of
this survey. \textbf{Resilience} emphasizes graceful degradation and recovery
\emph{after} a disruption (e.g., continued coordination once a fraction of
teammates fail). \textbf{Safety} constrains behavior to satisfy hard constraints
(costs, barrier conditions) and is robustness's frequent companion in
safety-critical applications. \textbf{Security} concerns an adversary with intent
and capabilities (an attack/threat model). These notions overlap---a certified
defense against communication attacks is simultaneously a robustness, security,
and (often) safety statement---and we flag the primary lens when a work spans
several.

Formally, robustness is defined relative to an \emph{uncertainty set}
$\mathcal{U}$ that collects the admissible models. Lifting robust dynamic
programming and robust MDPs \cite{nilim2005robust,iyengar2005robust} to the
multi-agent setting yields the \emph{robust Markov game} and its distributional
counterpart, the \emph{distributionally robust Markov game} (DRMG)
\cite{zhang2020,shi2024sample}, in which transitions (and sometimes rewards) may
lie anywhere in $\mathcal{U}$. The quantity of interest is the \emph{robust value}
\[
\underline{V}_i^{\boldsymbol{\pi}}(s)
\;=\; \inf_{P\in\mathcal{U}}\; V_i^{\boldsymbol{\pi}}(s; P),
\]
the return guaranteed against the worst admissible model, and the solution concept
becomes a \emph{robust equilibrium}---an equilibrium of the game in which nature
adversarially selects $P\in\mathcal{U}$. The shape of $\mathcal{U}$ (e.g.,
total-variation, Kullback--Leibler, or $R$-contamination balls; $(s,a)$- vs.\
$s$-rectangular sets) governs both tractability and conservativeness and recurs as
a design axis in \S\ref{sec:model}.

\paragraph{The curse of multiagency.} One MARL-specific obstacle deserves naming
up front because it reappears throughout. Because dynamics and rewards depend on
the \emph{joint} action, naive worst-case analysis scales with
$|\mathcal{A}|=\prod_{i}|\mathcal{A}_i|$, which is exponential in the number of
agents $n$. This \emph{curse of multiagency}---the explosion of the joint
action space---is the reason robust MARL is not a mechanical lift of single-agent
robust RL: guarantees, sample complexity, and equilibrium computation must all be
made to scale in $n$. Breaking this curse is the central theoretical thread of
\S\ref{sec:model} \cite{shi2024breaking,qu2026distributionally}.

\subsection{Threat/Uncertainty Model and the Survey Taxonomy}
\label{sec:bg-taxonomy}

Our organizing question is deliberately simple: \emph{what is being perturbed?}
Figure~\ref{fig:taxonomy} answers it. Branch~I gathers perturbation sources that
robust MARL \emph{inherits} from single-agent robust RL---uncertainty in the
environment/model (\S\ref{sec:model}), in states and observations
(\S\ref{sec:state}), and in actions and rewards (\S\ref{sec:adversarial}). These
are real and well studied, but they are not what makes the multi-agent problem
distinctive. Branch~II collects the dimensions that are \emph{absent in
single-agent RL and constitute the core argument of this survey}: attacked or
noisy \emph{communication} channels between agents (\S\ref{sec:comm}),
\emph{untrustworthy teammates}---Byzantine, faulty, or non-stationary co-players
(\S\ref{sec:byzantine}), and the \emph{curse of multiagency} on the theory side
(\S\ref{sec:model}). Branch~III marks emerging frontiers---LLM-based multi-agent
safety (\S\ref{sec:llm}) and offline/distribution-shift settings
(\S\ref{sec:offline})---and Branch~IV the cross-cutting concerns of benchmarks
(\S\ref{sec:benchmarks}) and applications (\S\ref{sec:applications}).

A \emph{secondary, methodological} axis cuts across all branches: distributionally
robust optimization, adversarial training, certified robustness (e.g., randomized
smoothing), and regularization are recurring \emph{tools} rather than perturbation
\emph{sources}. We therefore keep them off the primary axis and instead record,
per work, which tool is used in the comparison table of each section. A handful of
works span two perturbation sources (e.g., joint state-and-action adversarial
training \cite{bukharin2023robust}); we classify each in its primary branch and
cross-reference the secondary.

\subsection{Related Surveys and Our Differentiation}
\label{sec:bg-related}

Robustness has been surveyed thoroughly for the \emph{single-agent} case
\cite{moos2022robust,gu2026robust,mohan2026towards}, typically organized by
\emph{uncertainty source}---an axis we adopt and extend into the multi-agent
regime. General \emph{MARL} surveys \cite{zhang2021multi,gronauer2022multi,
jin2025comprehensive} chart algorithms and applications but treat robustness only
in passing. The closest prior work is the recent systematization of adversarial
machine learning in MARL \cite{adversarialml2025}, which is comprehensive on the
\emph{attack/defence} axis (state, communication, and partial Byzantine threats)
but, by construction, does not cover distributionally robust game \emph{theory}
(\S\ref{sec:model}), LLM multi-agent safety (\S\ref{sec:llm}), or
offline/distribution-shift robustness (\S\ref{sec:offline}).

Table~\ref{tab:related} contrasts coverage. Our contribution is not a new method
but a \emph{unifying frame}: we place attack/defence work, DRMG theory,
communication and Byzantine robustness, and the LLM/offline frontiers on a single
perturbation-source axis, and we use that frame to argue that the three Branch~II
dimensions---communication, untrustworthy teammates, and multiagency---are where
single-agent robustness tools do not transfer. These three dimensions recur as a
refrain in every subsequent section.

\begin{table}[t]\centering\footnotesize
\caption{Scope comparison with related surveys. \checkmark{} = covered as a
primary topic; \tcheck{} = partial/peripheral; -- = not covered.}
\label{tab:related}
\begin{tabular}{lccccccc}
\toprule
Survey & \makecell{Model\\unc.} & \makecell{State/\\Obs} & Comm. & Byzantine
 & \makecell{DRMG\\theory} & \makecell{LLM-\\MAS} & Offline \\
\midrule
Single-agent robust RL \cite{moos2022robust,gu2026robust} & \checkmark & \checkmark & -- & -- & \tcheck & -- & \tcheck \\
General MARL \cite{zhang2021multi,jin2025comprehensive} & -- & \tcheck & \tcheck & \tcheck & -- & \tcheck & -- \\
Adv.\ ML in MARL (SoK) \cite{adversarialml2025} & -- & \checkmark & \checkmark & \tcheck & -- & -- & -- \\
\textbf{This survey} & \checkmark & \checkmark & \checkmark & \checkmark & \checkmark & \checkmark & \checkmark \\
\bottomrule
\end{tabular}
\end{table}

\paragraph{Reading guide.} The next section opens with the deepest cluster, the
theory of robustness to environment/model uncertainty and the curse of
multiagency (\S\ref{sec:model}); \S\ref{sec:state}--\S\ref{sec:adversarial} treat
the inherited state/observation and action/reward dimensions;
\S\ref{sec:comm}--\S\ref{sec:offline} treat the MARL-specific and frontier
dimensions; and \S\ref{sec:benchmarks}--\S\ref{sec:applications} survey benchmarks and
applications before we distill open problems in \S\ref{sec:challenges}.

%=====================================================================
% §3 ENVIRONMENT / MODEL UNCERTAINTY (DRMG, theory)
% MASTER §1  (#1-27)        OWNER: Theory lead (Profile A)   |  TARGET: ~5 pp (deepest)
% DEPTH: deep -- state assumptions, uncertainty-set type, sample complexity, equilibrium.
% MUST CLOSE WITH:
%   (Q1) How is robustness to model uncertainty harder in MARL than single-agent?
%        -> curse of multiagency (joint action space explodes).
%   (Q2) Open gap? (function approx., online, infinite state, beyond TV/KL sets...)
%=====================================================================
\section{Robustness to Environment and Model Uncertainty}
\label{sec:model}

% Synthesize, do not enumerate. Organize into sub-themes:

\subsection{Formulations: Robust Markov Games and DRMGs}
% Minimax / robust MG \cite{} (Zhang NeurIPS'20 %TODO key); DRMG definition;
% uncertainty set families (TV, KL, $R$-contamination); mean-field type game \cite{zaman2024robust}.
TODO.

\subsection{Sample Complexity and the Curse of Multiagency}
% Core theory cluster: \cite{shi2024sample,shi2024breaking,qu2026distributionally,farhat2026sample}.
% Tabular vs function approximation (taming curses %TODO 2605.03125; strategically robust %TODO 2603.09208);
% online interaction \cite{farhat2026sample}; average-reward DRMG %TODO 2508.03136.
% Generative-model minimax %TODO NeurIPS'22; data-driven \cite{}; scalable \cite{}.
TODO.

\subsection{Risk-Sensitive / Risk-Averse Robustness}
% The 2025-26 convergence of robustness with risk aversion:
% \cite{mazumdar2024tractable,mazumdar2024behavioral,zhang2026provably}; mean-field risk %TODO 2602.13353.
TODO.

\subsection{Estimation, Evaluation, and Reward-Side Robustness}
% \cite{chen2025multi,li2026adaptive,kahe2026distributed}; correlated equilibrium %TODO;
% Bayesian distributional value %TODO; GOV-REK reward engineering %TODO; counterfactual %TODO.
TODO.

\begin{table}[t]\centering\footnotesize
\caption{Robust MARL under model/environment uncertainty (MASTER \S1).}
\label{tab:model}
\begin{tabular}{llllll}
\toprule
Work & Setting & Uncertainty set & Guarantee & Coop/Comp & Theory/Exp \\
\midrule
% one row per representative work; pull from robust-marl-MASTER.md §1
TODO & & & & & \\
\bottomrule
\end{tabular}
\end{table}

% [Closing paragraph] answer Q1 & Q2 explicitly -> tie to thesis (multiagency = MARL-specific).
   % MASTER §1
%=====================================================================
% §4 STATE / OBSERVATION PERTURBATION
% MASTER §2  (#28-42)       OWNER: Methods lead (Profile B)   |  TARGET: ~3 pp
% DEPTH: medium-high.  CLOSE WITH:
%   (Q1) Why harder in MARL? -> per-agent partial obs, critical-agent coupling,
%        joint worst-case state perturbation.
%   (Q2) Open gap? (scalable certificates, beyond $\ell_p$ perturbation, delayed/async obs)
%=====================================================================
\section{Robustness to State and Observation Perturbation}
\label{sec:state}

\subsection{Solution Concepts under State Adversaries}
% \cite{herobust,han2022solution} -- robust agent policy / worst-case state value.
TODO.

\subsection{Regularization- and Training-Based Defences}
% \cite{bukharin2023robust} (also adversarial, x-ref §5), \cite{li2023mir2} (MIR2),
% \cite{guo2024enhancing,guo2025robust}; RoMFAC %TODO; local advantage AC %TODO;
% temporal-consistency self-distillation \cite{chen2026enhancing}.
TODO.

\subsection{Partial Observability, Delay, and Agent Dropout}
% \cite{fu2025rainbow} (delayed obs), \cite{hazra2025tackling} (termination dynamics, x-ref §7),
% \cite{shi2025fault} (x-ref §7); local observations %TODO; adversaries-on-observation %TODO.
TODO.

\begin{table}[t]\centering\footnotesize
\caption{State/observation-robust MARL (MASTER \S2).}
\label{tab:state}
\begin{tabular}{lllll}
\toprule
Work & Perturbation model & Defence mechanism & Guarantee & Coop/Comp \\
\midrule
TODO & & & & \\
\bottomrule
\end{tabular}
\end{table}

% [Closing] answer Q1 & Q2.
   % MASTER §2
%=====================================================================
% §5 ADVERSARIAL ATTACKS & ADVERSARIAL TRAINING
% MASTER §3  (#43-59)       OWNER: Methods lead (Profile B)   |  TARGET: ~3.5 pp
% NOTE: this is the axis the CSUR'25 SoK covers -- be thorough AND show what we add.
% CLOSE WITH:
%   (Q1) Why harder in MARL? -> attacker can target critical agents / coordination;
%        single-agent poisoning can ruin the whole team.
%   (Q2) Open gap? (transferable defences, certified adversarial training at scale)
%=====================================================================
\section{Adversarial Attacks and Adversarial Training}
\label{sec:adversarial}

\subsection{Attack Vectors}
% Action: \cite{zhou2024adversarial}. State/critical-agent: Wolfpack \cite{lee2025wolfpack}.
% Black-box / partial obs: \cite{andam2025constrained,zhang2025black} (x-ref §10).
% Communication attacks: \cite{standen2026finding} (x-ref §6).
% Reward poisoning: reward-poisoning offline %TODO AAAI'23 (x-ref §9 offline).
% Application-domain attacks: \cite{alzubaidi2025adversarial} %TODO zhao action.
% Camouflage / stealth: \cite{lu2025camouflage}.
TODO.

\subsection{Adversarial Training and Robust Policy Learning}
% \cite{zhou2025robust} (stochastic adversary), interaction-breaking %TODO 2605.18024,
% \cite{hill2025co} (co-evolving curricula), \cite{peterson2025framework} (IsaacLab),
% M3DDPG %TODO (minimax), evolutionary auxiliary attackers %TODO,
% \cite{bukharin2023robust} (adversarial regularization, x-ref §4).
TODO.

\begin{table}[t]\centering\footnotesize
\caption{Attacks (top) and adversarial-training defences (bottom) for MARL (MASTER \S3).}
\label{tab:adversarial}
\begin{tabular}{lllll}
\toprule
Work & Attack/Defence & Vector (state/action/reward/comm) & Threat model & Coop/Comp \\
\midrule
TODO & & & & \\
\bottomrule
\end{tabular}
\end{table}

% [Closing] answer Q1 & Q2; explicitly contrast our coverage with attack/defence-only SoK.
         % MASTER §3
%=====================================================================
% §6 COMMUNICATION ROBUSTNESS  (MARL-specific dimension, Branch II)
% MASTER §4  (#60-72)       OWNER: Systems lead (Profile C)   |  TARGET: ~2.5 pp
% CLOSE WITH:
%   (Q1) Why MARL-specific? -> communication channel does not exist in single-agent RL.
%   (Q2) Open gap? (certified comm defences at scale, learned vs protocol comm)
%=====================================================================
\section{Robust and Secure Communication}
\label{sec:comm}

\subsection{Attacks on and Noise in Communication}
% malicious messages, perturbed channels; \cite{standen2026finding} (x-ref §5);
% noisy channels %TODO; mis-spoke/mis-lead %TODO.
TODO.

\subsection{Defences: Active Defence, Certification, Information-Theoretic}
% \cite{yu2024robust} (ADMAC active defence); certifiably robust comms %TODO (Furong Huang);
% certified policy smoothing %TODO; multi-view message certification %TODO;
% graph information bottleneck %TODO; temporal message control %TODO;
% \cite{ma2025robust} (decentralization-oriented adv. training); \cite{liu2026robust}.
TODO.

\subsection{Robust Communication in Applications}
% \cite{zhao2025towards} (multi-UAV, x-ref §11), \cite{xu2025dct} (platoon, x-ref §11);
% social empowerment %TODO.
TODO.

\begin{table}[t]\centering\footnotesize
\caption{Communication-robust MARL (MASTER \S4).}
\label{tab:comm}
\begin{tabular}{lllll}
\toprule
Work & Threat (noise/attack) & Defence & Certified? & Coop/Comp \\
\midrule
TODO & & & & \\
\bottomrule
\end{tabular}
\end{table}

% [Closing] answer Q1 & Q2 -- emphasize this is a dimension with NO single-agent analogue.
       % MASTER §4
%=====================================================================
% §7 UNTRUSTWORTHY TEAMMATES / BYZANTINE / FAULT TOLERANCE  (MARL-specific, Branch II)
% MASTER §5  (#73-86)       OWNER: Systems lead (Profile C)   |  TARGET: ~3 pp
% This + §6 are the STRONGEST evidence for the thesis (no single-agent analogue).
% CLOSE WITH:
%   (Q1) Why MARL-specific? -> "the adversary is a teammate"; trust/consensus problem.
%   (Q2) Open gap? (provable guarantees + scalability; detection vs robustness tradeoff)
%=====================================================================
\section{Untrustworthy Teammates, Byzantine Faults, and Trust}
\label{sec:byzantine}

\subsection{Byzantine-Robust Formulations}
% \cite{li2024byzantine} (Bayesian game), \cite{kazaribayesian} (unknown adversaries),
% \cite{mao2025ibgp} (IBGP, x-ref §8), fully Byzantine-resilient Q %TODO 2604.02791,
% decentralized Byzantine + reward machines %TODO.
TODO.

\subsection{Fault Tolerance and Agent Failure}
% \cite{mguni2025fault} (adversarial budget), \cite{shi2025towards} (fault tolerance),
% resilient value decomposition %TODO; robust multi-agent bandits %TODO (dishonest arms).
TODO.

\subsection{Trust- and Reputation-Based Mechanisms; Ad-hoc Teamwork}
% \cite{rudzitis2025trust} (UAV), \cite{hu2025trustorch} (x-ref §8), \cite{pan2025trust} (x-ref §11),
% \cite{patel2025modeling} (Werewolf trust/deception); unsupervised partner design (ad-hoc) %TODO.
TODO.

\begin{table}[t]\centering\footnotesize
\caption{Byzantine / fault-tolerant / trust-based robust MARL (MASTER \S5).}
\label{tab:byzantine}
\begin{tabular}{lllll}
\toprule
Work & Threat (Byzantine/failure/distrust) & Mechanism & Guarantee & Decentralized? \\
\midrule
TODO & & & & \\
\bottomrule
\end{tabular}
\end{table}

% [Closing] answer Q1 & Q2 -- this is the clearest "MARL adds a new dimension" case.
           % MASTER §5
%=====================================================================
% §8 ROBUSTNESS & SAFETY IN LLM MULTI-AGENT SYSTEMS  (frontier, differentiation highlight)
% MASTER §6  (#87-91)       OWNER: Systems/frontier lead (Profile C) or Lead | TARGET: ~2 pp
% This is the NEWEST sub-area and a key differentiator vs CSUR'25. Keep tight, forward-looking.
% CLOSE WITH:
%   (Q1) How does the LLM-agent setting reshape robust MARL? (NL channels, emergent roles)
%   (Q2) Open gap? (formal guarantees for LLM-MAS robustness; bridging RL-theory & LLM-agent practice)
%=====================================================================
\section{Robustness in LLM-Based Multi-Agent Systems}
\label{sec:llm}

\subsection{Internalized Safety via Co-Evolution}
% \cite{pan2025evo,pan2025advevo} (Evo-/AdvEvo-MARL, x-ref §5 adversarial co-evolution).
TODO.

\subsection{Constrained and Communication-Robust LLM Agents}
% \cite{wang2026safe} (natural-language constraints), \cite{hong2025collaborative} (edge LLM
% Stackelberg, x-ref §6), IBGP \cite{mao2025ibgp} (x-ref §7).
TODO.

\subsection{Position: Connecting Classical Robust MARL and LLM-Agent Security}
% \cite{} LLM-based MARL directions %TODO; relate to multi-agent security position papers.
% Argue the classical taxonomy (§3-§7) maps onto LLM-MAS threats.
TODO.

% [Closing] answer Q1 & Q2.
             % MASTER §6
%=====================================================================
% §9 OFFLINE / DISTRIBUTION-SHIFT ROBUSTNESS  (frontier)
% MASTER §7  (#92-93)       OWNER: Systems lead (Profile C)   |  TARGET: ~1 pp
% Short, emerging. CLOSE WITH:
%   (Q1) Why harder in MARL? -> joint-action OOD, exponential joint support.
%   (Q2) Open gap? (robust offline MARL with guarantees; connection to DRMG §3)
%=====================================================================
\section{Offline and Distribution-Shift Robustness}
\label{sec:offline}

% \cite{jin2026partial} (partial action replacement), \cite{li2025sample} (robust offline
% self-play, x-ref §3 model-based). Relate distribution shift to DRMG uncertainty sets (§3).
% Reward-poisoning on offline MARL %TODO (x-ref §5).
TODO.

% [Closing] answer Q1 & Q2; note this bridges to §3 (a robust-offline + DRMG synthesis is open).
             % MASTER §7
%=====================================================================
% §10 BENCHMARKS & EVALUATION  (cross-cutting)
% MASTER §8  (#94-105)      OWNER: Systems lead (Profile C) / entry (Profile D) | TARGET: ~2 pp
% CLOSE WITH:
%   (Q1) Are current benchmarks adequate to measure MARL-specific robustness?
%   (Q2) Open gap? (standardized robustness metrics; agent-failure & communication-attack benches)
%=====================================================================
\section{Benchmarks and Evaluation}
\label{sec:benchmarks}

\subsection{Robustness-Oriented Benchmarks and Testbeds}
% \cite{} Robust Gymnasium %TODO; \cite{li2025starcraft+} (adversary paradigm);
% \cite{barta2025measuring} (partial agent failure, x-ref §7); general MARL: SMACv2 %TODO, MATE %TODO, Mava %TODO.
TODO.

\subsection{Robustness Testing and Empirical Studies}
% \cite{} robustness testing on critical agents %TODO (x-ref §2,§5); comprehensive testing %TODO;
% empirical study on robustness & resilience %TODO 2510.11824 (x-ref §7);
% observation poisoning %TODO; domain studies (traffic signal, urban driving, x-ref §11).
TODO.

\begin{table}[t]\centering\footnotesize
\caption{Benchmarks/testbeds for robust MARL (MASTER \S8).}
\label{tab:bench}
\begin{tabular}{lllll}
\toprule
Benchmark & Perturbation types & Metric & Agents/scale & Open-source \\
\midrule
TODO & & & & \\
\bottomrule
\end{tabular}
\end{table}

% [Closing] answer Q1 & Q2 -- motivates a "standardized robustness benchmark" challenge in §12.
          % MASTER §8
%=====================================================================
% §11 APPLICATIONS  (cross-cutting)
% MASTER §9  (#106-146)     OWNER: Lead + entry (Profile D)   |  TARGET: ~3 pp
% DEPTH: shallow/compact. Do NOT write up all 41 papers individually.
% Format: one short paragraph per domain + ONE big summary table. Pick representatives.
% CLOSE WITH:
%   (Q1) Which robustness dimensions does each application domain stress most?
%   (Q2) Open gap? (sim-to-real robustness, deployment-grade guarantees)
%=====================================================================
\section{Applications}
\label{sec:applications}

% One compact paragraph per domain; cite representatives, defer rest to Table~\ref{tab:apps}.
\subsection{Connected/Autonomous Vehicles \& Safe-Robust Control}
% \cite{zhang2025safety,smith2025robust,dong2025robustness,wang2025robust} %TODO miao 5G; sim-to-real.
\subsection{Power Grids and Energy}
% \cite{pu2025distributed,tian2025robust,shah2025uncertainty,ghasemi2024combating} %TODO + [M] grid apps.
\subsection{Traffic, Routing, and Logistics}
% \cite{pei2025distributionally,vieira2025decentralized,low2026robust,liu2025distributionally} (x-ref §3 DRMG).
\subsection{UAV/USV, Robotics, Networks, and Cyber-Physical Security}
% \cite{li2025romuc,khan2025robust,amarnath2026robust,shen2025rmaac,thomas2026multi,holz2025towards,alqithami2025hierarchical}.
\subsection{Finance and Other Domains}
% \cite{chen2026mars}.
TODO.

\begin{table*}[t]\centering\footnotesize
\caption{Robust MARL applications by domain (MASTER \S9, representatives).}
\label{tab:apps}
\begin{tabular}{lllll}
\toprule
Domain & Work & Robustness dimension stressed (\S3--\S9) & Robust technique & Real/Sim \\
\midrule
CAV/AV       & TODO & state / comm / safety & & \\
Power/energy & TODO & model unc. / adversarial & & \\
Traffic/logistics & TODO & DRMG / dist. shift & & \\
UAV/robot/cyber   & TODO & comm / Byzantine & & \\
Finance      & TODO & risk-averse & & \\
\bottomrule
\end{tabular}
\end{table*}

% [Closing] answer Q1 & Q2.
        % MASTER §9
%=====================================================================
% §12 OPEN CHALLENGES & FUTURE DIRECTIONS   OWNER: Lead (synthesize from all sections)
% TARGET: ~3 pp  -- THE PAYOFF SECTION; reviewers weigh this heavily.
% Each challenge = 1 short subsection: state problem, why open, point to enabling work.
%=====================================================================
\section{Open Challenges and Future Directions}
\label{sec:challenges}

\subsection{Unifying the Three MARL-Specific Dimensions}
% No framework jointly handles model unc. + untrustworthy teammates + comm attacks.
\subsection{Breaking the Curse of Multiagency Beyond Tabular}
% function approx / online / infinite state are nascent (\S3) -- guarantees still partial.
\subsection{Scalable Certified Robustness for Cooperative Teams}
% certification (\S4,\S6) does not scale with agent count.
\subsection{From Attacks to Provable Defences}
% attack literature (\S5) outpaces defences with guarantees.
\subsection{Robustness in LLM Multi-Agent Systems}
% bridge classical robust-MARL theory (\S3-\S7) with LLM-agent practice (\S8).
\subsection{Standardized Robustness Benchmarks and Metrics}
% \S10 shows fragmentation; need agent-failure + comm-attack standardized suites.
\subsection{Sim-to-Real and Deployment-Grade Guarantees}
% \S11 apps mostly simulation; deployment robustness is open.
\subsection{Risk-Sensitivity, Offline, and Distributional Robustness Unification}
% \S3 risk-averse + \S9 offline + DRMG are converging but not unified.
TODO prose for each.

%=====================================================================
% §13 CONCLUSION                         OWNER: Lead   |  TARGET: ~1 pp
%=====================================================================
\section{Conclusion}
\label{sec:conclusion}

% Recap: unified perturbation-source taxonomy over 146 works (2017-2026).
% Restate the thesis payoff: MARL inherits single-agent robustness dimensions AND
% adds untrustworthy teammates / attacked communication / curse of multiagency.
% Point to §12 as the agenda. End with the "national-interest" framing (safe deployment
% of multi-agent autonomy) -- one crisp closing sentence.
TODO.


\bibliographystyle{ACM-Reference-Format}
\bibliography{refs}
% refs.bib: build from raw-data/rmarl-paper-from-25-to-26.bib AFTER de-duplication
% (he2023robust x5, shi2024breaking x2, zaman2024robust x2, qu=learning..., etc.)

\end{document}
